\documentclass[12pt]{extarticle}
\usepackage[a4paper,left=2.8cm,right=2.8cm,top=2.2cm,bottom=2cm]{geometry}
\usepackage{amsmath,xcolor,enumitem,amssymb}
\pagestyle{empty}
\setlength{\parindent}{0pt}
\setlist{leftmargin=1.4em,itemsep=3pt,topsep=4pt}
\newcommand{\abschnitt}[1]{\par\vspace{18pt}{\large\bfseries #1}\par\vspace{5pt}}
\newcommand{\jj}{\,$\bullet$}
\begin{document}
{\LARGE\bfseries Abitur · Mathematik GK}\hfill{\small\color{gray}Stand Oktober 2026}
\abschnitt{Termin}
Mittwoch, 5.\,Mai 2027, 9 Uhr, 285 Minuten (Berlin und Brandenburg)
\abschnitt{Zwei Teile}
\begin{tabular}{@{\hspace{1.4em}}l@{\hspace{1em}}l@{}}
\textbf{Teil A} & kurze Aufgaben \textbf{ohne} Taschenrechner und Formelsammlung,\\
 & Abgabe spätestens nach 100 Minuten\\[3pt]
\textbf{Teil B} & große Aufgaben \textbf{mit} Hilfsmitteln, Analysis mit Wahl\\
\end{tabular}
\abschnitt{Mitbringen}
Taschenrechner (wie im Unterricht: WTR oder MMS/CAS), Formelsammlung des IQB, Geodreieck
\abschnitt{In dieser Reihenfolge lernen}
{\small Anteil an allen Punkten der Prüfungen 2022--2026. Jedes Thema kam in fast jedem Jahr dran.}\par\vspace{6pt}
\begin{minipage}{12cm}
\textbf{Analysis} \hfill \textbf{55\,\%}\par
\begin{enumerate}[label=\arabic*.,leftmargin=1.8em,itemsep=2pt]
\item Ableitung + Tangente \hfill 13\,\%
\item Kurvenuntersuchung \hfill 29\,\%
\item Integral \hfill 13\,\%
\end{enumerate}\vspace{4pt}
\textbf{Stochastik} \hfill \textbf{23\,\%}\par
\begin{enumerate}[label=\arabic*.,leftmargin=1.8em,itemsep=2pt]
\item Baumdiagramm + bedingte Wahrscheinlichkeit \hfill 13\,\%
\item Binomialverteilung \hfill 6\,\%
\item Erwartungswert \hfill 3\,\%
\end{enumerate}\vspace{4pt}
\textbf{Geometrie} \hfill \textbf{23\,\%}\par
\begin{enumerate}[label=\arabic*.,leftmargin=1.8em,itemsep=2pt]
\item Punkte, Flächen, Körper \hfill 9\,\%
\item Geraden + Ebenen \hfill 6\,\%
\item Winkel + Abstände \hfill 7\,\%
\end{enumerate}\vspace{4pt}
\end{minipage}
\par\vspace{8pt}
{\small Die Hälfte aller Punkte kommt aus der Analysis. Ihr Kern -- Nullstellen, Extrem- und Wendepunkte -- kommt jedes Jahr; vor der Prüfung gezielt wiederholen.}
\abschnitt{Alte Prüfungen zum Üben}
Stark-Band Abitur 2027 (Grundkurs): Prüfungen 2022--2025 mit Lösungen, 2026 online.
\end{document}
